\documentclass[11pt]{article}

\usepackage[margin=1in]{geometry}
\usepackage{amsmath,amssymb,bm}
\usepackage{booktabs}
\usepackage{graphicx}
\usepackage{subfig}
\usepackage{times}
\usepackage{url}

\newcommand{\calD}{\mathcal{D}}
\newcommand{\Nfive}{N_{5\%}}
\newcommand{\distdyn}{d^{\mathrm{dyn}}}
\newcommand{\figurepdf}[2]{%
    \IfFileExists{#2}{\includegraphics[width=#1]{#2}}{%
        \fbox{\parbox[c][3.4cm][c]{#1}{\centering\small
        \textbf{PDF FIGURE PENDING}\\[1.5mm]\texttt{#2}}}}%
}

\renewcommand{\thefigure}{S\arabic{figure}}
\renewcommand{\thetable}{S\arabic{table}}
\renewcommand{\theequation}{S\arabic{equation}}

\title{Supplementary Material: Fleet-Informed Bayesian Optimization for Data-Efficient Controller Calibration}
\author{Jakob Weber \and Markus Gurtner \and Adrian Trachte \and Andreas Kugi}
\date{}

\begin{document}
\maketitle

This document contains implementation and robustness details that are intentionally kept outside the L4DC main submission to preserve the ten-page main-text budget. The main paper is self-contained with respect to its claims and conclusions.

\section{Detailed LQI Parameterization}

For yaw-rate tracking, the identified two-state model uses $\mathbf{x}=[\beta,r]^\top$ and is augmented by the integral tracking error
\begin{equation}
 z_{k+1}=z_k+T_s(r_k^{\rm ref}-r_k).
\end{equation}
With $\mathbf{x}_k^a=[\mathbf{x}_k^\top,z_k]^\top$, the fixed scale vector is
\begin{equation}
 \mathbf{s}=[0.053229,\,0.575573,\,0.139862]^\top,
\end{equation}
and the LQI penalties are
\begin{equation}
 \mathbf Q(\boldsymbol{\xi})
 =\operatorname{diag}\!\left(
 \frac{10^{\xi_1}}{s_1^2},
 \frac{10^{\xi_2}}{s_2^2},
 \frac{10^{\xi_3}}{s_3^2}\right),
 \qquad \mathbf R=1.
\end{equation}
The scales were fixed before the confirmatory BO study from representative nominal closed-loop trajectories and were not tuned per vehicle family.

\section{Bayesian-Optimization Implementation}

\begin{table}[h]
\centering
\caption{BO implementation used in the held-out study.}
\begin{tabular}{@{}ll@{}}
\toprule
Quantity & Setting\\
\midrule
Search box & $[-3,3]^3$\\
Local initialization & 2 Sobol points\\
Acquisition & Expected improvement, minimization\\
EI exploration margin & $0.01$\\
EI candidate set & 256 Sobol points\\
Output scaling & $J_i/J_i^{\rm base}$\\
Source weighting floor & $0.02$\\
GP kernel & Constant $\times$ Mat\'ern-$5/2$\\
GP length scale & $\boldsymbol{\ell}_0=0.3\mathbf{1}$, bounds $[0.05,3]^3$\\
GP base noise $\alpha_0$ & $10^{-4}$\\
GP optimizer restarts & 2\\
Infeasible penalty & $3.0$ in normalized units\\
Baseline $\boldsymbol{\xi}^{\rm base}$ & $[0,\,0.30103,\,0.30103]^\top$\\
\bottomrule
\end{tabular}
\end{table}

At target evaluation $n$, expected improvement is maximized over the fixed 256-point Sobol candidate set. For dynamics-weighted transfer, source-record noise is scaled as
\begin{equation}
 \alpha_{ij}=\frac{\alpha_0}{\max(s_{ij},0.02)}.
\end{equation}

\section{Vehicle Model, Fleet Generation, and Identification}

The continuous-time linear control-design model is
\begin{equation}
 \dot{\mathbf{x}}=\mathbf A\mathbf x+\mathbf B\delta,
\end{equation}
with
\begin{equation}
\mathbf{A}=\begin{bmatrix}
-\dfrac{C_f+C_r}{mU} & \dfrac{-aC_f+bC_r}{mU^2}-1\\[2mm]
\dfrac{-aC_f+bC_r}{I_z} & -\dfrac{a^2C_f+b^2C_r}{I_zU}
\end{bmatrix},
\qquad
\mathbf{B}=\begin{bmatrix}\dfrac{C_f}{mU}\\[2mm]\dfrac{aC_f}{I_z}\end{bmatrix}.
\end{equation}
The model is discretized by zero-order hold with $T_s=0.05$\,s.

Controller candidates are evaluated on a nonlinear lateral plant. The slip angles are
\begin{equation}
 \alpha_f=\delta-\beta-\frac{ar}{U},
 \qquad
 \alpha_r=-\beta+\frac{br}{U},
\end{equation}
and the tire forces are
\begin{equation}
 F_{y,q}=\mu F_{z,q}\tanh\!\left(\frac{C_q\alpha_q}{\mu F_{z,q}}\right),
 \qquad q\in\{f,r\},
\end{equation}
with static axle loads $F_{z,f}=mgb/(a+b)$ and $F_{z,r}=mga/(a+b)$. The nonlinear states evolve as
\begin{equation}
 \dot\beta=\frac{F_{y,f}+F_{y,r}}{mU}-r,
 \qquad
 \dot r=\frac{aF_{y,f}-bF_{y,r}}{I_z}.
\end{equation}

\begin{table}[h]
\centering
\caption{Vehicle-family prototypes and common simulation parameters.}
\begin{tabular}{@{}lccc@{}}
\toprule
Parameter & Nominal & Payload & Tire-degraded\\
\midrule
$m$ [kg] & 1500 & 1800 & 1500\\
$I_z$ [$\mathrm{kg\,m^2}$] & 2700 & 3375 & 2700\\
$a$ [m] & 1.25 & 1.25 & 1.25\\
$b$ [m] & 1.35 & 1.35 & 1.35\\
$C_f$ [$\mathrm{N/rad}$] & 80000 & 80000 & 60000\\
$C_r$ [$\mathrm{N/rad}$] & 90000 & 90000 & 90000\\
$U$ [m/s] & 15 & 15 & 15\\
$T_s$ [s] & 0.05 & 0.05 & 0.05\\
$\mu$ [-] & 1.0 & 1.0 & 1.0\\
\bottomrule
\end{tabular}
\end{table}

For each generated system, $m$, $I_z$, $C_f$, $C_r$, and $U$ are independently multiplied by $1+\eta$, $\eta\sim\mathcal U[-0.025,0.025]$, while $a$ and $b$ remain fixed. A system is resampled if the fixed baseline controller is infeasible on the calibration task, with at most 50 attempts.

The local discrete model is estimated as
\begin{equation}
 \boldsymbol{\Theta}_i=[\mathbf A_{d,i}\;\mathbf B_{d,i}]\in\mathbb R^{2\times3},
 \qquad
 \hat{\boldsymbol{\theta}}_i=\operatorname{vec}(\hat{\boldsymbol{\Theta}}_i)\in\mathbb R^6,
\end{equation}
using matrix RLS with regressor $[\mathbf x_k^\top,u_k]^\top$, nominal-model initialization, forgetting factor $\lambda=0.99$, and initial covariance $\mathbf P_0=\mathbf I_3$. Each system processes eight repeated 6\,s identification episodes. If $\mathbf D_i=\operatorname{diag}(\hat\sigma_{i,1}^2,\hat\sigma_{i,2}^2)$ contains the final row-wise innovation variances and $\mathbf P_i$ is the regressor covariance, the precision used for dynamics distance is
\begin{equation}
 \hat{\mathbf W}_i=\operatorname{pinv}\!\left(\mathbf D_i\otimes\mathbf P_i\right).
\end{equation}

\section{Calibration Task and Held-Out Protocol}

Every candidate is evaluated on three stochastic episodes for each of two lane-change half-durations, $T_{\rm lc}\in\{1,2\}$\,s. For a trajectory with $T$ samples,
\begin{equation}
 J_i(\boldsymbol{\xi},\omega)
 =\sqrt{\frac{1}{T}\sum_{k=1}^{T}
 \left(r_{i,k}(\boldsymbol{\xi},\omega)-r_k^{\rm ref}\right)^2}.
\end{equation}
The unfiltered reference is zero until $T_0=1$\,s, follows a positive sine lobe of duration $T_{\rm lc}$ and then a negative sine lobe of equal duration, with peak $r_{\max}=35^\circ/$s. A zero-phase first-order Butterworth filter with 1\,Hz cutoff smooths the waveform. Each episode lasts 6\,s. Feasibility requires successful LQI synthesis, $|\beta|\le6^\circ$, and steering-rate change no larger than $0.2618$\,rad per simulation step.

\begin{figure}[h]
\centering
\figurepdf{0.62\textwidth}{figures/fig6_lane_change_reference.pdf}
\caption{Yaw-rate references used for controller calibration.}
\end{figure}

\begin{table}[h]
\centering
\caption{Final held-out evaluation protocol.}
\begin{tabular}{@{}ll@{}}
\toprule
Quantity & Setting\\
\midrule
Source systems per realization & 30, 10 per family\\
Held-out fleet realizations & 20\\
Targets per realization & 6, 2 per family\\
Source history per system & 25 independent-BO evaluations\\
Warm-start history per target & 725 source observations\\
Local Sobol initialization & 2 evaluations\\
Maximum target budget & 12 evaluations\\
Methods & independent / global / dynamics-weighted\\
Primary metric & restricted mean $\Nfive$\\
\bottomrule
\end{tabular}
\end{table}

For the source-maturity sensitivity study, each source history is generated once from two Sobol initialization evaluations followed by sequential independent BO to 25 total evaluations. The nested prefixes are $M_{\rm hist}\in\{2,5,10,25\}$. The main mature-fleet comparison instead uses eight source initialization evaluations followed by 17 sequential BO evaluations, so the two 25-evaluation endpoints need not be numerically identical.

\section{Additional Result Diagnostics}

\begin{figure}[h]
\centering
\subfloat[Cumulative observed target attainment.]{\figurepdf{0.47\textwidth}{figures/fig4a_n5pct_ecdf.pdf}}
\hfill
\subfloat[Fixed-budget success probabilities.]{\figurepdf{0.47\textwidth}{figures/fig4b_n5pct_success_probability.pdf}}
\caption{Additional threshold-based views of target-system sample efficiency. Unresolved independent-BO runs remain right-censored rather than being treated as successes at the final budget.}
\end{figure}

The source-coverage experiment holds the source pool at nine systems and changes the number of same-family peers. It tests whether the global-versus-weighted conclusion depends on full fleet access or on a particular source composition.

\begin{figure}[h]
\centering
\figurepdf{0.62\textwidth}{figures/fig5b_coverage_sensitivity.pdf}
\caption{Sensitivity to same-family source coverage in a nine-peer source pool.}
\end{figure}

The second source-maturity view reports the same paired result as experiment savings rather than absolute restricted mean.

\begin{figure}[h]
\centering
\figurepdf{0.62\textwidth}{figures/fig7b_fleet_maturity_savings.pdf}
\caption{Paired mean number of target evaluations saved relative to independent BO versus source maturity.}
\end{figure}

\end{document}
